\documentclass[12pt]{article}
\usepackage{booktabs}
\usepackage{inputenc, amsmath,amssymb}
\usepackage[font=small,labelfont=bf,justification=justified,format=plain]{caption}
\usepackage{subcaption}
\usepackage{graphicx}
\usepackage{xcolor}
\usepackage{float}
\usepackage[colorlinks=true,linkcolor=blue,citecolor=blue,urlcolor=blue]{hyperref}
\usepackage{array}
\usepackage{natbib}
\usepackage{setspace}
\usepackage[margin=.75in]{geometry}
\usepackage{amsmath, amssymb, amsfonts}
\usepackage{threeparttable}

\onehalfspacing

\graphicspath{{../../../../../Dropbox/Lawsuits Asymmetry/output}}

\title{Entitlement Pipeline: Environmental Review Congestion --- Data Exploration}
\author{Tyler Jacobson}

\begin{document}
\maketitle

\section{Overview}

This document reports on the structure of the environmental review dataset (\texttt{env\_post\_2015}) as a precursor to testing whether congestion in the environmental review process affects project completion. Environmental review cases are filed under the prefix \texttt{ENV} and are distinct from the main entitlement case. The CEQA type determines the level of scrutiny: Categorical Exemptions (CE) require no further analysis, while Environmental Impact Reports (EIR) involve full-scale review.

The total sample contains 14,630 post-2015 environmental review cases. The central question for the congestion analysis is whether variation in queue length at the time of filing affects how long the environmental review takes, and whether that in turn affects whether the underlying entitlement project completes.

\section{CEQA Type Distribution}

\begin{figure}[H]
  \centering
  \includegraphics[width=0.6\textwidth]{congestion/env/bar_ceqa_type_share.pdf}
  \caption{Share of environmental review cases by CEQA type.}
\end{figure}

Categorical Exemptions account for the large majority of cases (12,062 out of 14,630, or 82\%). The remaining cases involve genuine environmental analysis: MND (1,147 cases), EAF (966 cases), EIR (233 cases), ND (135 cases), and a small unclassified residual (87 cases).

The dominance of CE cases is consequential for the congestion analysis. CE determinations are administrative: the planner declares the project exempt from CEQA at the time of filing, so there is little scope for a queue to affect their length. \textbf{The relevant sample for congestion analysis is the non-CE cases} (MND, ND, EIR), where a genuine review process occurs.

\section{Case Length by CEQA Type}

Table~\ref{tab:stats_by_type} reports case length and completion rates by CEQA type. Several patterns stand out.

\begin{table}[H]
\centering
\begin{threeparttable}
\caption{Environmental Review Case Statistics by CEQA Type}
\label{tab:stats_by_type}
\begin{tabular}{lrrrrrrr}
\toprule
Type & N & Mean (days) & Median & P90 & \% Completed & \% Failed \\
\midrule
CE           & 12,062 &  74 &   0 &  251 & 99.8 &  3.3 \\
ND           &    135 & 274 & 143 &  724 & 93.3 &  1.5 \\
MND          &  1,147 & 265 & 175 &  531 & 87.9 &  3.8 \\
EAF          &    966 & 622 & 456 & 1436 & 41.9 & 56.3 \\
EIR          &    233 & 997 & 900 & 2121 & 26.2 &  1.3 \\
Unclassified &     87 & 429 & 302 &  929 & 54.0 &  4.6 \\
\bottomrule
\end{tabular}
\begin{tablenotes}
  \small
  \item Notes: Length is in calendar days from filing to completion. Cases with negative-length records are excluded. \% Failed = share with determination coded as denied, dismissed, terminated, or withdrawn.
\end{tablenotes}
\end{threeparttable}
\end{table}

\textbf{CE cases:} The median CE length is 0 days, meaning the majority complete on the day they are filed. This is consistent with how CE works in practice --- the planner marks the project exempt at intake, generating an instantaneous determination. These cases will be excluded from the main congestion regressions.

\textbf{MND and ND cases:} Median lengths of 175 and 143 days, respectively. These are the cases most likely to show congestion effects, as they involve real analytical work but are not so complex as to be dominated by project-specific factors.

\textbf{EIR cases:} A median of 900 days (2.5 years) and only 26\% completion. The low completion rate primarily reflects right-censoring: EIRs filed close to the end of the sample window have not had time to complete. The distribution of EIR lengths will be heavily censored and will require a hazard approach.

\textbf{EAF cases:} The 56\% failure rate is anomalous and warrants investigation. EAF (Environmental Assessment Form) cases may be a transitional or reclassified track rather than a completed review type. It is worth checking whether EAF cases that ``fail'' are subsequently refiled under a different CEQA type.

\begin{figure}[H]
  \centering
  \begin{subfigure}[t]{0.48\textwidth}
    \centering
    \includegraphics[width=\textwidth]{congestion/env/hist_env_length.pdf}
    \caption{All ENV cases}
  \end{subfigure}
  \hfill
  \begin{subfigure}[t]{0.48\textwidth}
    \centering
    \includegraphics[width=\textwidth]{congestion/env/kdensity_env_length_by_type.pdf}
    \caption{By CEQA type (CE excluded)}
  \end{subfigure}
  \caption{Distribution of environmental review case length.}
\end{figure}

The overall histogram is dominated by CE cases at zero. The KDE plot shows the length distributions for the substantive review types. EIR cases have a tight cluster around 900 days with a heavy right tail. MND and ND cases overlap considerably in the 100--500 day range, which will limit power if they are pooled.

\section{Time Series of Filings}

\begin{figure}[H]
  \centering
  \begin{subfigure}[t]{0.48\textwidth}
    \centering
    \includegraphics[width=\textwidth]{congestion/env/ts_env_filings.pdf}
    \caption{All ENV filings}
  \end{subfigure}
  \hfill
  \begin{subfigure}[t]{0.48\textwidth}
    \centering
    \includegraphics[width=\textwidth]{congestion/env/ts_env_filings_by_type.pdf}
    \caption{By CEQA type}
  \end{subfigure}
  \caption{Monthly environmental review case filings.}
\end{figure}

The overall volume time series shows any structural breaks in filing rates, including the expected drop during the COVID-19 period (2020--2021). The type-stratified panel is the key diagnostic for the congestion analysis: the MND queue should show meaningful time-series variation to identify the congestion effect.

\section{Right-Censoring}

\begin{figure}[H]
  \centering
  \includegraphics[width=0.6\textwidth]{congestion/env/bar_censored_by_year.pdf}
  \caption{Share of ENV cases uncompleted, by filing year.}
\end{figure}

Right-censoring rises mechanically for more recent filing cohorts. Cases filed in 2022--2023 have had little time to complete. For the hazard model, this is handled naturally by the Cox likelihood. The concerning pattern to watch for is elevated censoring in \textit{pre-2020} cohorts, which would suggest projects stalled rather than completed or officially withdrawn.

\section{Staff Assignment Analysis}

\begin{figure}[H]
  \centering
  \begin{subfigure}[t]{0.48\textwidth}
    \centering
    \includegraphics[width=\textwidth]{congestion/env/hist_gap_assigned_env.pdf}
    \caption{Days from filing to staff assignment}
  \end{subfigure}
  \hfill
  \begin{subfigure}[t]{0.48\textwidth}
    \centering
    \includegraphics[width=\textwidth]{congestion/env/hist_cases_per_reviewer_env.pdf}
    \caption{Total cases per reviewer}
  \end{subfigure}
  \caption{Staff assignment characteristics for ENV cases.}
\end{figure}

The gap between filing and staff assignment is the key data quality diagnostic. As noted in the data cleaning paper, the assignment date for some ENV cases reflects an extension review rather than the original assignment (the Laura Steele / William Hughen example). Cases with assignment dates that come very late (e.g., more than a year after filing) or \textit{before} the filing date are likely to have unreliable staff attribution and should be flagged or excluded from the reviewer-level analysis.

The cases-per-reviewer distribution determines whether there are enough cases per reviewer to identify workload effects. A highly skewed distribution where a few reviewers handle most cases is helpful for power; a flat distribution where every reviewer only handles a handful of cases would limit the reviewer-level design.

\begin{figure}[H]
  \centering
  \includegraphics[width=0.6\textwidth]{congestion/env/scatter_reviewer_eir_share.pdf}
  \caption{Reviewer specialization: share of EIR cases vs.\ total caseload (reviewers with $\geq 10$ cases).}
\end{figure}

Reviewers who specialize in EIR cases will have very different workload profiles than those who handle primarily CE and MND. This specialization needs to be accounted for in the reviewer-level workload regressions (Step 6 of the plan) --- either by stratifying by CEQA type or by using predicted case length as the workload measure rather than raw case counts.

\section{Implications for Congestion Analysis}

\subsection{Sample restrictions}

The congestion analysis should restrict to non-CE cases (MND, ND, EIR). CE cases complete instantly and will not show congestion effects; including them would dilute the queue measure and add noise. The working sample is approximately 2,500 cases.

\subsection{Defining the queue}

The queue measure should be defined separately for substantive review types. An MND reviewer's congestion comes from other MND/ND cases, not from EIR cases (which require a different level of effort). The next step (Step 2) builds this stratified queue.

\subsection{EAF cases}

The anomalously high failure rate for EAF cases (56\%) needs to be understood before including them. If EAF represents a preliminary screening that routes cases to MND/EIR, then the ``failure'' is not a true project failure but a reclassification. Investigate whether failed EAF cases have linked MND or EIR cases filed shortly after.

\subsection{Staff reliability}

Run the hazard analysis both with and without the reviewer-level controls. If the assignment date diagnostic shows a large share of late or pre-filing assignments, the reviewer-level design should use only cases with assignment dates that fall within a reasonable window after filing (e.g., 0--180 days).

\section{Reviewer-Level MND/ND Queue}

\subsection{Queue Construction}

The queue for reviewer $r$ at the start of month $m$ is the number of MND/ND cases assigned to $r$ that were filed before month $m$ and had not yet completed by the start of month $m$. Formally, using flow accounting:
\[
  Q_{rm} = Q_{r,m-1} + \text{Filed}_{r,m-1} - \text{Finished}_{r,m-1},
  \qquad Q_{r,m_0} = 0,
\]
where $m_0$ is the first month in the sample. The initialization at zero understates the true queue early in the sample because cases filed before 2015 are not observed. The first 12 months of each reviewer's queue are therefore excluded from all downstream regressions.

Cases with assignment dates that precede the filing date or that arrive more than 365 days after filing are dropped as unreliable (consistent with the Laura Steele assignment issue documented in the data cleaning paper).

\subsection{Queue Distribution}

\begin{figure}[H]
  \centering
  \begin{subfigure}[t]{0.48\textwidth}
    \centering
    \includegraphics[width=\textwidth]{congestion/env/hist_reviewer_queue_mnd_nd.pdf}
    \caption{Histogram of queue at filing}
  \end{subfigure}
  \hfill
  \begin{subfigure}[t]{0.48\textwidth}
    \centering
    \includegraphics[width=\textwidth]{congestion/env/ts_reviewer_queue_mnd_nd.pdf}
    \caption{Queue distribution over time}
  \end{subfigure}
  \caption{Reviewer-level MND/ND queue. The histogram shows the number of active cases in a reviewer's queue at the time each focal case is filed (trimmed at the 99th percentile). The time series shows the monthly cross-sectional distribution of queue lengths across reviewers with at least one active case; the shaded band is the interquartile range.}
\end{figure}

The histogram describes the identifying variation available for the hazard regression: cases that arrived when their reviewer had few other pending cases versus those that arrived during a backlog. The time series reveals whether queue lengths have systematic trends or seasonality that should be absorbed by time fixed effects in the regression.

\section{Hazard Analysis: Does Queue Length Slow Environmental Reviews?}

\subsection{Design}

The hazard regression asks whether MND/ND cases filed when their reviewer's queue is longer take longer to complete. The estimating equation is a Cox proportional hazards model:
\[
  h_j(t) = h_{0r(j)}(t) \exp\!\bigl(\delta \log(Q_{r(j),m(j)} + 1) + X_j'\beta\bigr),
\]
where $h_{0r}(t)$ is a reviewer-specific baseline hazard, $Q_{r,m}$ is the reviewer's queue at the start of the filing month, and $X_j$ includes a MND indicator (versus ND). A hazard ratio below 1 on $\log(Q+1)$ means a longer queue reduces the instantaneous rate of completion---cases take longer.

Column (1) uses no fixed effects and clusters standard errors by reviewer. Column (2) adds filing-year fixed effects to absorb secular changes in processing speed. Column (3) uses a stratified Cox model with reviewer-specific baseline hazards, identifying solely off within-reviewer variation in queue length over time; this is the most conservative specification and closest to a causal estimate.

\subsection{Results}

{\footnotesize
\input{../../../../../Dropbox/Lawsuits Asymmetry/output/congestion/env/hazard_queue_mnd_nd.tex}
}

\subsection{Binscatter}

\begin{figure}[H]
  \centering
  \begin{subfigure}[t]{0.48\textwidth}
    \centering
    \includegraphics[width=\textwidth]{congestion/env/binscatter_env_length_queue.pdf}
    \caption{Raw correlation}
  \end{subfigure}
  \hfill
  \begin{subfigure}[t]{0.48\textwidth}
    \centering
    \includegraphics[width=\textwidth]{congestion/env/binscatter_env_length_queue_fes.pdf}
    \caption{Controlling for reviewer and year-month FEs}
  \end{subfigure}
  \caption{ENV review length and reviewer queue at filing. The right panel nets out reviewer and year-month fixed effects, isolating within-reviewer variation in queue from secular trends.}
\end{figure}

The raw binscatter shows the unconditional relationship between queue length and review duration. The panel with reviewer and year-month fixed effects is the most informative: the slope here represents the within-reviewer congestion effect, purged of both reviewer quality differences and aggregate time trends.

\end{document}
